\newtheorem{theorem}{Theorem}[section] 
\newtheorem{satz}[theorem]{Satz} 
\newtheorem{definition}[theorem]{Definition} 
\newtheorem{lemma}[theorem]{Lemma} 
\newtheorem{beispiel}[theorem]{Beispiel} 
\newtheorem{alg}[theorem]{Algorithmus} 

\newcommand{\ket}[1]{\left| #1 \right>}
\newcommand{\proj}[2]{|#1\left.\right>\left<\right.#2|} 
\newcommand{\scalar}[2]{\left<#1|#2\right>} 
\newcommand{\scalars}[3]{\left<#1|#2|#3\right>} 
\newcommand{\reg}[2]{\ket{\Psi\left( #1,#2\right)}} 
\newcommand{\dft}{\mbox{DFT}} 
\newcommand{\diag}[1]{\mbox{ diag}\left(#1\right)} 
\newcommand{\unten}[1]{\lfloor #1 \rfloor} 
\newcommand{\oben}[1]{\lceil #1 \rceil} 
\newcommand{\betrag}[1]{\vert #1 \vert} 
\newcommand{\k}{\log_{\lambda} m_0} 
\newcommand{\ok}{\lceil \log_{\lambda} m_0 \rceil} 

\newenvironment{beweis}{{\bf Beweis: }}{\hfill *}
 
\Vortrag{Quantenalgorithmen}{Pawel Wocjan}
\noindent 

%% *************************************** Einfuhrung ******************** 
 \section{Einf"uhrung} 
   Quantencomputer bieten die M"oglichkeit auf koh"arent "uberlagerten Zust"anden parallel zu rechnen und die Interferenz zwischen verschiedenen Berechnungspfaden 
   auszunutzen. Der "ubliche Weg, das Vorhandensein 
   eines Zustands festzustellen, ist die Durchf"uhrung einer Messung, also einer Projektion auf orthogonale Eigenr"aume einer geeigneten Observablen. Die 
   Wahrscheinlichkeit, den gew"unschten Zustand zu messen, ist das Betragsquadrat seiner Amplitude in der "Uberlagerung bez"uglich der durch 
   die Eigenvektoren der Observablen definierten Orthonormalbasis. Ungl"ucklicherweise kann es jedoch passieren, da\3 das Betragsquadrat der Amplitude des gesuchten Zustands
   exponentiell klein ist. In einem solchen Fall sind durchschnittlich exponentiell viele Messungen notwendig, 
   um den gew"unschten Zustand festzustellen. Die einfachste Vorgehensweise, in der resultierenden "Uberlagerung die gew"unschte Ausgabe festzustellen, 
   scheitert also an den physikalischen Prinzipien der Messung eines quantenmechanischen Systems. 
  
   Im folgenden werden drei Algorithmen behandelt, die durch eine iterative Anwendung einer unit"aren Operation die Amplituden der L"osungen verst"arken und 
   damit die Wahrscheinlichkeit, eine L"osung bei der Messung zu erhalten, erheblich erh"ohen. Leider ben"otigt dieser Verst"arkungsproze\3 eine 
   exponentielle Laufzeit. Trotzdem sind diese Algorithmen schneller als alle m"oglichen klassischen Algorithmen.  
%% ************************* Grovers Algorthmus ************************** 
 \section{Grovers Suchalgorithmus} 
%% ************************ Einfuhrung *********************************** 
  \subsection{Problemstellung} 
    Sei $X_N=\left\{0,1,\ldots ,N-1\right\}$ f"ur ein $N\in \N$ und sei $C:X_N \to \left\{ 0,1\right\}$ eine beliebige 
    Funktion, die angibt, ob $i\in X_N$ eine bestimmte Bedingung $C$ erf"ullt . Das Ziel ist es, ein solches $i\in X_N$ zu finden, 
    so da\3 $C\left(i\right)=1$ gilt, vorausgesetzt, da\3 so ein $i$ existiert. Falls $C$ berechnet werden kann, man aber absolut keine
    Information "uber $i\in X_N$ mit $C(i)=1$ aus der Berechnungsvorschrift f"ur $C$ herausfinden kann, das hei\3t, wenn $C(i)=0$ oder $C(i)=1$ nur durch
    Anwendung von $C$ entschieden werden kann, dann kann kein klassischer Algorithmus (deterministisch oder probabilistisch) eine Erfolgswahrscheinlichkeit 
    von mehr als $50\%$ erreichen, ohne $C$ f"ur mehr als $N/2$ $i\in X_N$ auszuwerten. Jeder klassische 
    Algorithmus braucht also ${\cal O}(N)$ Schritte, um diese Aufgabe zu l"osen. Grover \cite{grosuch} entdeckte einen Algorithmus f"ur  
    einen Quantencomputer, der dieses Problem in einer Zeit ${\cal O}(\sqrt{ N})$ l"osen kann, vorausgesetzt, da\3 eine  
    eindeutige L"osung existiert. Grovers Algorithmus l"a\3t sich f"ur mehrere L"osungen verallgemeinern, wobei die Anzahl der L"osungen $t$ vorher aber
    bekannt sein mu\3. Am Ende des Algorithmus wird eine beliebige L"osung mit einer bestimmten Wahrscheinlichkeit geliefert. Die Anzahl der L"osungen ist
    dabei entscheidend; ansonsten funktioniert der Algorithmus nicht. Er hat eine Laufzeit von ${\cal O}(\sqrt{N/t})$. Der Algorithmus l"a\3t sich weiter
    so modifizieren, da\3 er auch bei einer unbekannten Anzahl von L"osungen am Ende der Berechnung eine beliebige L"osung liefert. Dieser Algorithmus
    hat eine Laufzeit von ${\cal O}(\sqrt{N/t})$. Diese Algorithmen wurden in \cite{hoyer} vorgestellt.
    Anzahl der L"osungen zu sch"atzen.

%% ************************* Beschreibung des Algorithmus ****************  
  \subsection{Beschreibung des Algorithmus} 
    Grovers Algorithmus besteht aus einer Initialisierung, gefolgt von einer bestimmten Anzahl $m$ identischer Schritte, 
    einer Messung und einem klassischen Test. F"ur jede Funktion $C:X_N \to \left\{ 0,1\right\}$ sei die bedingte  
    Phasen\-drehung $S_C$ definiert durch 
    $$ S_C \ket{i} = \left\{ 
                        \begin{array}{cc}
                              \ket{i} & \mbox{ falls $C\left( i\right) = 1$} \\ 
                              \ket{i} & \mbox{ sonst }
                        \end{array}
                     \right.
    $$
    In einer praktischen Implementation der bedingten Phasendrehung $S_C$ mu\3 der Inhalt des Quantenregisters, der die Argumente aufnimmt, festgestellt 
    werden, und es mu\3 entschieden werden, ob dessen Phase um $\pi$ gedreht werden soll oder nicht. Es darf nach dieser Operation nichts in den
    Hilfsregistern, die zur Ausf"uhrung dieser Operation ben"otigt werden, "ubrigbleiben, um die Interferenzen nicht zu beeinflu\3en. 
    \begin{enumerate} 
      \item Der Ausgangspunkt dieser Operation ist $\ket{i,0,{\bf 0}}$.
        Dabei ist $0$ eine freie Speicherstelle zur Aufnahme von $C(i)$ sowie ${\bf 0}$ ein hinreichend gro\3er Zwischenspeicher zur Berechnung  
        von $C$. 
      \item Berechne $C(i)$ durch die unit"are Transformation $U_C$ f"ur alle "uberlagerten Zust"ande $\ket{i}$ 
        $$ U_C \left( \sum_{i=0}^{N-1} \alpha_i \ket{i,0,{\bf 0}} \right) = \sum_{i=0}^{N-1} \alpha_i \ket{i,C(i),\Psi_i}. $$ 
        Dabei bezeichnet $\ket{\Psi_i}$ einen nicht n"aher spezifizierten Zustand, der w"ahrend der Berechnung von $C(i)$ entsteht. 
      \item F"uhre eine Phasenverschiebung um $(-1)^{C(i)}$ durch. Dies liefert den Zustand 
        $$ \sum_{i=0}^{N-1} (-1)^{C(i)} \alpha_i \ket{i,C(i),\Psi_i} $$ 
      \item Stelle f"ur jedes i den Zustand $\ket{i,0,{\bf 0}}$ wieder her durch R"uckg"angigmachen der Berechnung von $C$ durch 
        $$ U_C^{-1} \left( \sum_{i=0}^{N-1} (-1)^{C(i)} \alpha_i \ket{i,C(i),\Psi_i} \right) =  
           \sum_{i=0}^{N-1}(-1)^{C(i)} \alpha_i \ket{i,0,{\bf 0}} $$ 
           Dieser Schritt hei\3t "`uncomputing"'. Er erm"oglicht die ungest"orte Interferenz der Zust"ande $\ket{i}$. 
    \end{enumerate} 
    Der wichtigste Unterschied zwischen Quantenmechanik und klassischer Mechanik ist der, da\3 ein Quantenzustand des Gesamtsystems im allgemeinen nicht in die Zust"ande
    von Untersystemen zerlegt werden kann. Deswegen d"urfen nach einer Quantenberechnung keine Zwischenergebnisse "ubrigbleiben, weil sonst die erhaltenen Ergebnisse nicht
    in einer f"ur die Quantencomputer notwendigen Art und Weise verwendet werden k"onnte. Nach der bedingten Phasendrehung $S_C$ mu\3 $\ket{\Psi_i}$
    durch "`uncomputing"' r"uckg"angig gemacht werden. F"ur eine genaue Erkl"arung siehe \cite{kitaev}.
    
    In jedem Iterationsschritt sind also zwei Auswertungen der Funktion $C$ notwendig.
    Weiterhin braucht man f"ur Grovers Algorithmus eine lokale unit"are Transformation $W$, die zus"atzlich die Eigenschaft, aus dem Zustand $\ket{0}$ die Gleichverteilung 
    $$ W\ket{0} = \frac{1}{\sqrt{N}} \sum_{i=0}^{N-1} \ket{i} $$  
    herzustellen, erf"ullt. Diese Transformation mu\3 lokal sein, damit sie physikalisch implementiert werden kann. 
    Algebraisch bedeutet dies, da\3 sie sich als Tensorprodukt von kleinen Matrizen schreiben l"a\3t. 
    Falls $N=2^n$ eine Zweierpotenz ist, dann ist die Walsh-Hadamard-Transformation (die allgemeine Fouriertransformation $\dft_{\oplus_{i=1}^{n} Z_2}$ 
    zu der Gruppe $\oplus_{i=1}^{n} Z_2$) eine Transformation, die sich physikalisch sehr einfach realisieren l"a\3t, da auf jeden Qubit  
    lediglich eine "`fair coin"'-Operation ausgef"uhrt wird, und die die gew"unschte Eigenschaft besitzt. Falls $N$ keine Zweierpotenz ist, kann man die
    Fouriertransformation $\dft_{Z_N}$ zu einer zyklischen Gruppe $Z_N$ der Ordnung $N$ nehmen. Diese Transformation l"a\3t sich auf einem 
    Quantencomputer realisieren.
    
    Die Anwendung der Operation $W$ auf den Zustand $\ket{0}$ liefert die "Uberlagerung aller $n$-Qubit Eingaben mit gleicher Amplitude und  
    Phase. Diese Konfiguration der Gleichverteilung der Zust"ande ist der Ausgangspunkt vieler Quantenalgorithmen, da nun alle m"oglichen Eingaben in der  
    "Uberlagerung zur parallelen Abarbeitung bereitstehen. 
    
    Eine weitere Transformation, die ben"otigt wird, ist die {\em Diffusionsoperation} $D$, die wie folgt definiert ist 
    $$ D = W R W^{\dagger}, $$ 
    wobei $R = \mbox{diag}\{1,-1,\ldots,-1\}$. Die {\em Grovertransformation} $G_C$ zu der Bedingung $C$ ist die unit"are Transformation, die durch
    die Hintereinanderausf"uhrung der bedingten Phasendrehung $S_C$ und der Diffusionstransformation $D$ 
    $$ G_C = D S_C $$ 
    definiert ist.
   
    Grovers Algorithmus stellt zun"achst den Gleichverteilungszustand $\ket{\Psi} = W\ket{0}$ her. Dann wird $G_C$ $m$-mal nacheinander 
    angewandt, wobei $m$ noch bestimmt werden mu\3. Schlie\3lich wird der Zustand gemessen. 
    Der Algorithmus hat nur dann Erfolg, wenn sich nach der Messung der Zustand $\ket{i}$ ergibt und $C(i)=1$ gilt. 
%% ********************************** Interpretation der Diffusionstransformation ************************* 
    \subsection{Interpretation der Diffusionstransformation} 
      Die Diffusionstransformation $D$ kann als eine Inversion am Durchschnittswert betrachtet werden. Eine einfache Inversion ist eine 
      Phasenrotation. Sei $\alpha$ der Durchschnittswert der Amplituden aller Zust"ande $\ket{i}$ . Nach Anwendung der  
      Diffusionstransformation wird die Amplitude jedes Zustands derart ver"andert, da\3 sie jetzt um genauso viel unterhalb ("uber)  
      $\alpha$ liegt, wie sie vor der Transformation "uber (unterhalb) $\alpha$ lag. 
      \begin{satz} \label{diffusion}
        Die Diffusionstransformation sieht folgenderma\3en aus 
        $$ D = \pmatrix{ -1+\frac{2}{N} & \frac{2}{N}   & \ldots      & \frac{2}{N}   \cr 
                               \frac{2}{N}   & \ddots        & \ddots      & \vdots        \cr 
                               \vdots        & \ddots        & \ddots      & \frac{2}{N}   \cr 
                               \frac{2}{N}   & \ldots        & \frac{2}{N} & -1+\frac{2}{N} \cr}. $$ 
      \end{satz}   
      \beweis   
        Die Diffusionstransformation $D$ ist durch $ D = W R W^{\dagger}$ definiert. W ist eine beliebige lokale unit"are Transformation, die 
        $ W\ket{0} = \frac{1}{\sqrt{N}} \sum_{i=0}^{N-1} \ket{i} $ erf"ullt. Damit sehen die Transformation $W$ und ihre Adjugierte 
        $W^{\dagger}$ folgenderma\3en aus 
        $$ W = \frac{1}{\sqrt{N}}\pmatrix{1      & *     & \ldots & *      \cr 
                                          \vdots & \vdots& \ddots & \vdots \cr 
                                          1      & *     & \ldots & *      \cr} $$ 
        $$ W^{\dagger} = \frac{1}{\sqrt{N}}\pmatrix{ 1      & \ldots & 1      \cr 
                                                     *      & \ldots & *      \cr  
                                                     \vdots & \ddots & \vdots \cr 
                                                     *      & \ldots & *      \cr}. $$ 
        Um zu sehen, da\3 $D$ diese Gestalt hat, mu\3 man $D$ als Summe von zwei einfacheren Matrizen schreiben 
        \begin{eqnarray*} 
          D & = & W \diag{1,-1,\ldots ,-1} W^{\dagger} \\ 
            & = & W \diag{-1,-1,\ldots ,-1} W^{\dagger} + W \diag{2,0,\ldots, 0} W^{\dagger} \\ 
            & = & -I + \frac{2}{\sqrt{N}}\pmatrix{1      & 0     & \ldots & 0      \cr 
                                          \vdots & \vdots& \ddots & \vdots \cr 
                                          1      & 0     & \ldots & 0      \cr} W^{\dagger} \\ 
            & = & -I + \frac{2}{N}\pmatrix{1      & \ldots & 1      \cr 
                                          \vdots & \ddots  & \vdots \cr 
                                          1      & \ldots & 1 \cr} \\ 
            & = & -I + 2P.
         \end{eqnarray*} 
          $D$ kann als Summe $D=-I+2P$ geschrieben werden, wobei $I$ die Einheitsmatrix und $P$ eine Matrix ist, f"ur deren Koeffizienten 
          $P_{ij}=\frac{1}{N}$ $ (i,j=0,\ldots N-1)$ gilt. $P$ sieht folgenderma\3en aus
          $$ P = \frac{1}{N} \sum_{i,j=0}^{N-1} \proj{i}{j}. $$ 
          Die Anwendung von $P$ auf einen beliebigen Zustand $\ket{\Psi}$ ergibt einen Zustand, dessen Koeffizienten gleich dem Durchschnittswert $\alpha$ der  
          Komponenten von $\ket{\Psi}$ sind. 
          \begin{eqnarray*}       
            P\ket{\Psi} & = & \left( \frac{1}{N} \sum_{i,j=0}^{N-1} \proj{i}{j} \right) \sum_{k=0}^{N-1} \scalar{k}{\Psi}\ket{k} \\ 
                        & = & \frac{1}{N} \sum_{i,j,k=0}^{N-1} \scalar{k}{\Psi} \ket{i}\scalar{j}{k} \\ 
                        & = & \frac{1}{N} \sum_{i,j=0}^{N-1} \scalar{j}{\Psi}\ket{i} \\ 
                        & = & \left( \frac{1}{N} \sum_{j=0}^{N-1}\scalar{j}{\Psi} \right) \sum_{i=0}^{N-1} \ket{i} \\ 
                        & = & \sum_{i=0}^{N-1} \alpha \ket{i} 
          \end{eqnarray*} 
          Um zu sehen, da\3 $D$ eine Inversion am Durchschnittswert realisiert, mu\3 man betrachten, wie $D$ auf einem beliebigen Zustand $\ket{\Psi}$ 
          wirkt. 
          \begin{eqnarray*} 
            D\ket{\Psi} & = & \left( -I+2P\right) \ket{\Psi} \\ 
                        & = & -\ket{\Psi} + 2P\ket{\Psi} 
          \end{eqnarray*} 
          F"ur die $i$-te Komponente dieses Zustands gilt 
          \begin{eqnarray*} 
            \scalars{i}{D}{\Psi} & = & -\scalar{i}{\Psi} + 2\alpha \\  
                                 & = & \left( \alpha + \left( \alpha - \scalar{i}{\Psi}\right)\right) . 
          \end{eqnarray*} 
          \hfill $\Box$
          Dies entspricht genau der Inversion am Durchschnittswert. In einer Groveriteration werden durch die bedingte Phasendrehung $S_C$ zun"achst  
          die Amplituden der L"osungen mit $-1$ multipliziert und liegen danach unter dem Durchschnittswert $\alpha$. Nach Anwendung der  
          Diffusionstransformation sind sie positiv und gr"o\3er als zuvor. Man sieht also, da\3 mit jeder Iteration die Amplituden der L"osungen  
          wachsen. Allerdings ist dies nicht unbegrenzt der Fall, denn wenn die Amplituden der L"osungen zu gro\3 werden, wird der Durchschnitt nach  
          der bedingten Phasendrehung $S_C$ negativ. Damit fallen die Amplituden der L"osungen zugunsten der anderen Zust"ande ab. Es kommt also sehr 
          genau darauf an, die Anzahl der Iterationen zu bestimmen. Eine genaue Kenntnis des geeigneten Me\3zeitpunktes ist f"ur eine  
          praktische Anwendung unerl"a\3lich. 

%% **************************** Bekannte Anzahl der Losungen ********************************** 
       \subsection{Bekannte Anzahl der L"osungen} 
         Betrachten wir nun den Fall, wenn $t$ L"osungen existieren, das hei\3t, es gibt $t$ verschiedene Werte f"ur $i\in X_N$ mit $C(i)=1$. Wir sind an einer 
         beliebigen L"osung interessiert. Sei die Anzahl der L"osungen $t$ als bekannt vorausgesetzt. Definiere $A=\{ i|C(i)=1 \}$ und  
         $B=\{ i|C(i)=0 \}$. 
         F"ur zwei reelle positive Zahlen $k$ und $l$ mit $tk^2+(N-t)l^2=1$ definiere den Zustand des Quantenregisters durch 
         $$ \reg{k}{l} = \sum_{i\in A} k\ket{i} + \sum_{i\in B} l\ket{i} .$$ 
         Die Anwendung einer Groveriteration transformiert die Amplituden $k$ und $l$ folgenderma\3en 
         \begin{eqnarray*} 
            k_{j+1} & = & \frac{N-2t}{N}k_j + \frac{2(N-t)}{N}l_j \\ 
            l_{j+1} & = & \frac{N-2t}{N}l_j - \frac{2t}{N}k_j . 
         \end{eqnarray*} 
         Um den optimalen Zeitpunkt f"ur die Messung zu bestimmen mu\3 man diese Rekurrenz l"osen und die Amplituden $k_j$ und $l_j$ in einer 
         geschlossenen Form darstellen. Diese Rekurrenzrelation l"a\3t sich als Matrixmultiplikation 
         $$ \pmatrix{k_{j+1} \cr l_{j+1}} = \pmatrix{ \frac{N-2t}{N} & \frac{2(N-t)}{N} \cr 
                                                      -\frac{2t}{N}  & \frac{N-2t}{N} } 
                                            \pmatrix{k_j \cr l_j} $$
         schreiben. 
         F"ur die $j$-te Iteration des Groversalgorithmus gilt 
         $$ \pmatrix{k_{j} \cr l_{j}} = \pmatrix{ \frac{N-2t}{N} & \frac{2(N-t)}{N} \cr 
                                                 -\frac{2t}{N}   & \frac{N-2t}{N} }^j 
                                            \pmatrix{1/\sqrt{N} \cr 1/\sqrt{N}}. $$ 
         Die Anfangsbedingungen $k_0=l_0=1/\sqrt{N}$ ergeben sich aus der Gleichverteilung, die am Anfang hergestellt wird.
         Die Rekurrenzrelation l"a\3t sich mit dem folgenden Mathematica-Programm l"osen.
         \begin{center}{\small 
         \begin{verbatim}
(* Grover-Iteration, SE, PW, 18.12.96, Mathematica v2.0 *)

(* Setze tN = t/N = Sin[theta]^2 mit 0 < theta < Pi/2 *)
G = {
  { 1 - 2 tN, 2 (1 - tN) },
  { -2 tN,    1 - 2 tN   }
};

(* spezielle Vereinfachung *)
simplify[X_] :=
  Expand[
    X //. {
      Sqrt[Sin[theta]^2]      -> Sin[theta],
      1/Sqrt[Sin[theta]^2]    -> 1/Sin[theta],
      Sqrt[-1 + Sin[theta]^2] -> I Cos[theta]
    },
    Trig -> True
  ]

(* G diagonalisieren: G = A . DiagonalMatrix[d] . Inverse[A] *)
ES = Eigensystem[G] /. tN -> Sin[theta]^2; 
d  = simplify[ ES[[1]] ];
A  = simplify[ Transpose[ES[[2]]] ];

(* ==> die Eigenwerte von G sind 
   d = { Cos[2 theta] - I Sin[2 theta], 
         Cos[2 theta] + I Sin[2 theta]
       } 
*)

(* Damit ist die geschlossene Form: *)
g = 
  ComplexExpand[
    A . 
    DiagonalMatrix[{Exp[-2 I j theta], Exp[2 I j theta]}] . 
    Inverse[A]
  ];



(* ==> 
   g = { { Cos[2 j theta],               Cot[theta] Sin[2 j theta] },
         { -(Sin[2 j theta] Tan[theta]), Cos[2 j theta]            }
       }
*)

(* Mit den Anfangsbedingungen der Rekurrenz
     k[0] = l[0] = 1/Sqrt[N] und
     {k[j+1], l[j+1]} = G . {k[j], l[j]}
   ist die geschlossene L"osung
*)
kl = g . {1/Sqrt[N], 1/Sqrt[N]};

(* Damit ergibt sich die Loesung fuer die Amplituden {k_j, l_j} =
         Cos[2 j theta]   Cot[theta] Sin[2 j theta]
        {-------------- + -------------------------, 
            Sqrt[N]                Sqrt[N]
 
         Cos[2 j theta]   Sin[2 j theta] Tan[theta]
         -------------- - -------------------------}
            Sqrt[N]                Sqrt[N]
*)
         \end{verbatim} } \end{center}
         Unter Verwendung der Additionstheoreme f"ur trigonometrische Funktionen bekommt man folgende geschlossene Formeln f"ur die Amplituden der L"osungen und
         der Nicht-L"osungen 
         \begin{eqnarray*} 
           k_{j} & = & \frac{1}{\sqrt{t}}\sin\left(\left( 2j + 1\right)\Theta\right) \\ 
           l_{j} & = & \frac{1}{\sqrt{N-t}}\cos\left(\left( 2j + 1\right)\Theta\right) . 
         \end{eqnarray*} 
         Der Winkel $\Theta$ ist dabei so gew"ahlt, da\3 $\sin^2\Theta =\frac{t}{N}$ gilt. Es gilt $\Theta \in \left(0,\frac{\pi}{2} \right)$. 
         Die Erfolgswahrscheinlichkeit ist am gr"o\3ten, wenn $\betrag{l_m}^2$ minimal ist.
         $$ l_{\tilde m} = 0 \iff \tilde m = \left( \pi - 2\Theta \right) / 4\Theta + k\cdot \frac{\pi}{2\Theta}, $$ 
         wobei $k\in\Z.$ Im folgenden wird $k=0$ gew"ahlt.
         Da die Anzahl der Iterationen $m$ eine nat"urliche Zahl sein mu\3, wird $m$ folgenderma\3en definiert 
         $$ m = \unten{\pi / 4\Theta}.$$ 
         Der optimale Zeitpunkt f"ur eine Messung w"are nach $ \tilde m = \left( \pi - 2\Theta \right) / 4\Theta $ Iterationen. Es wird gezeigt, da\3 nach
         $m$ Iterationen die Erfolgswahrscheinlichkeit sehr gut ist. \\ 
         \begin{tabular}{crcl} 
                         & $\betrag{m - \tilde m}$ &   $=$   & $\betrag{\unten{\pi /4\Theta} - \pi /4\Theta + 1/2} \le 1/2$          \\ 
           $\Rightarrow$ & $\betrag{\left( 2\tilde m + 1\right)\Theta - \left( 2m + 1\right)\Theta}$ & $\le$ & $\Theta$              \\ 
           $\Rightarrow$ & $\betrag{\cos\left(\left(2m+1\right)\Theta\right)}$                       & $\le$ & $\betrag{\sin\Theta}$ \\ 
         \end{tabular} \\ 
         Die Wahrscheinlichkeit, da\3 ein falscher Wert nach $m$ Iterationen gemessen wird, ist gleich 
         $$ \left( N-t\right) l_m^2 = \cos^2\left(\left(2m+1\right)\Theta\right) \le \sin^2\Theta = t/N. $$ 
         Sie ist vernachl"a\3igbar, falls $t \ll N$. Dieser Algorithmus hat eine Laufzeit von ${\cal O}\left(\sqrt{\frac{N}{t}}\right)$, da 
         $\Theta \ge \sin\Theta = \sqrt{t/N}$ und damit  
         $$ m \le \frac{\pi}{4\Theta} \le \frac{\pi}{4} \sqrt{\frac{N}{t}} $$ 
         gilt. 
        
         Bemerkung: 
         \begin{enumerate}
           \item Eventuell liefert $k\neq 0$ ein besseres Ergebnis.
           \item F"uhre eine gebrochene Iteration $G^{1/\nu}$ f"ur ein geeignetes $\nu\in\Z$, um besser an die Nullstelle heranzukommen.
         \end{enumerate}

%% ********************** Untere Schranke ************************
       \subsection{Untere Schranke f"ur die Laufzeit von Quantensuchalgorithmen}
         Grover weist daraufhin, da\3 jeder Quantensuchalgorithmus mindestens eine zu $\sqrt{N}$ proportionale Laufzeit hat, wenn es eine eindeutige L"osung gibt, wobei er
         sich auf allgemeine Aussagen aus \cite{bernstein} st"utzt. Diese Aussage l"a\3t sich soweit verallgemeinern, da\3 man eine explizite untere Grenze f"ur die
         Anzahl der Auswertung der Funktion $C$ als Funktion der Anzahl der L"osungen $t$ angeben kann. 
         \begin{satz}
           Sei $S$ eine beliebige Menge von $N$ Strings, und $M$ eine bounded-error Quantenturingmaschine, die mit einem Orakel $O$ ausgestattet ist. Sei $A\subset S$ eine
           zuf"allig ausgew"ahlte Untermenge von $S$ der M"achtigkeit $t$ mit $t\ge 1$, wobei alle diese Untermengen gleichwahrscheinlich sind. Das Orakel $O$ sei so definiert, 
           da\3 $O(x)=1$ genau dann wenn $x\in A$. Die erwartete Anzahl, wie oft $M$ das Orakel $O$ befragen mu\3, um ein beliebiges $y\in A $ mit einer Wahrscheinlichkeit von
           mindestens $1/2$ zu finden, ist mindestens $\unten{(\sin(\pi/8))\sqrt{\unten{N/t}}}.$
         \end{satz}
         \beweis siehe \cite{bernstein} \hfill $\Box$

         Grovers Suchalgorithmus braucht $m\le\unten{(\pi/4)\sqrt{N/t}}$ Iterationen. Da in jedem Iterationsschritt, die bedingte Phasendrehung $S_C$ ben"otig wird, werden also
         insgesamt mindestens $2\unten{\frac{\pi}{4}\sqrt{N/t}}$ Auswertungen der Funktion $C$ ausgef"uhrt. 

%% ************************* Unbekannte Anzahl der Losungen ******************************************* 
       \subsection{Unbekannte Anzahl der L"osungen} 
         Die Anzahl der Iterationen $m$ h"angt entscheidend von der Anzahl der L"osungen $t$ ab. Grovers Algorithmus kann also nicht direkt verwendet  
         werden, falls die Anzahl der L"osungen unbekannt ist. Der Algorithmus mu\3 geeignet modifiziert werden. Es werden die folgenden zwei Lemmata
         bei der Analyse des Algorithmus ben"otigt.
         \begin{lemma} F"ur eine ganze positive Zahl $m$ und eine beliebige reelle Zahl $\alpha$ mit $\alpha \neq k\pi, k\in\Z$ gilt
           $$ \sum_{j=0}^{m-1} \cos\left(\left( 2j+1\right) \alpha\right) = \frac{\sin\left( 2m\alpha\right)}{2\sin \alpha}. $$ 
         \end{lemma} 
         \beweis Durch vollst"andige Induktion nach $m$ \hfill $\Box$
         \begin{lemma} 
           Sei $t$ die Anzahl der unbekannten L"osungen. Sei der Winkel $\Theta$ so definiert, da\3 $\sin^2\Theta = t/N$ gilt. Sei $m$ eine positive  
           ganze Zahl. Die Zahl $j$ werde zuf"allig aus $\{ 0,\ldots, m-1\}$ ausgew"ahlt, wobei alle Zahlen gleichwahrscheinlich seien. Die Wahrscheinlichkeit, eine 
           L"osung nach $j$ Groveriterationen zu erhalten, ist genau 
           $$ P_m = \frac{1}{2} - \frac{\sin(4m\Theta)}{4m\sin(2\Theta)}. $$ 
           Insbesondere ist $P_m \ge \frac{1}{4}$ falls $m \ge \frac{1}{\sin(2\Theta)}$. 
         \end{lemma} 
         \beweis 
           Die Erfolgswahrscheinlichkeit nach $j$ Groveriterationen ist gerade $tk_j^2 = \sin^2((2j+1)\Theta)$. Die durchschnittliche  
           Erfolgswahrscheinlichkeit, wenn j aus $\{0,\ldots,m-1\}$ zuf"allig gew"ahlt wird, ist 
           \begin{eqnarray*} 
              P_m & = & \frac{1}{m} \sum_{j=0}^{m-1} \sin^2((2j+1)\Theta) \\
                  & = & \frac{1}{2m} \sum_{j=0}^{m-1} \left( 1-\cos((2j+1)2\Theta)\right) \\ 
                  & = & \frac{1}{2} - \frac{1}{2m} \sum_{j=0}^{m-1} \cos((2j+1)2\Theta) \\ 
                  & = & \frac{1}{2} - \frac{\sin(4m\Theta)}{4m\sin(2\Theta)}.  
           \end{eqnarray*} 
           Es wird die Identit"at $\sin^2(x)=\frac{1}{2}(1-\cos(2x))$ verwendet.
           Falls $m \ge \frac{1}{\sin(2\Theta)}$ ist, dann gilt $\frac{\sin(4m\Theta)}{4m\sin(2\Theta)} \le \frac{1}{4m\sin(2\Theta)} \le \frac{1}{4}$ und  
           damit $P_m \ge \frac{1}{4}$. \hfill $\Box$

           Wir k"onnen jetzt einen Suchalgorithmus angeben, der eine beliebige L"osung findet, auch wenn die Anzahl $t$ aller L"osungen unbekannt ist.
           \subsection*{Algorithmus}
             Es wird angenommen, da\3 $1 \le t \le \frac{3}{4}N$ gilt. 
             \begin{enumerate} 
               \item Setze $m = 1$ und $\lambda = \frac{6}{5}$. 
               \item W"ahle $j$ zuf"allig aus den nichtnegativen ganzen Zahlen kleiner als $m$, wobei jede dieser Zahlen gleichwahrscheinlich ist. 
               \item F"uhre $j$ Iterationen des Groveralgorithmus beginnend im Zustand $\ket{\Psi_0} = \frac{1}{\sqrt{N}} \sum_{i_0}^{N-1} \ket{i}$ aus.
               \item Messe das Quantenregister. Sei $i$ das Ergebnis 
               \item Falls $C(i)=1$, dann wurde eine L"osung gefunden, stop.
               \item Sonst setze $m = \min(\lambda m, \sqrt{N})$ und gehe zu 2. 
             \end{enumerate} 
           \begin{satz} 
             Dieser Algorithmus findet eine L"osung in ${\cal O}(\sqrt{N/t})$. 
           \end{satz} 
           \beweis 
              Sei $\Theta$ so, da\3 $\sin^2\Theta = t/N$ gilt. Setze $m_0 = \frac{1}{\sin(2\Theta)}$. Dann gilt 
              \begin{eqnarray*} 
                m_0 & = & \frac{1}{\sin(2\Theta)} \\ 
                    & = & \frac{1}{2\sin(\Theta)\cos(\Theta)} \\
                    & = & \frac{1}{2\sqrt{t/N}\sqrt{(N-t)/t}} \\
                    & = & \frac{1}{2}\sqrt{\frac{N}{t}}\sqrt{\frac{t}{N-t}} \\
                    & = & \frac{N}{2\sqrt{(N-t)N}} \\ 
                    & < & \sqrt{N/t}
              \end{eqnarray*} 
              f"ur $t<\frac{3}{4}N < \frac{4}{5}N$.
              \hfill $\Box$

              Wir werden die erwartete Anzahl der ausgef"uhrten Groveriterationen absch"atzen. Bei dem $s$-ten Durchlauf der Schleife gilt  
              $m=\lambda^{s-1}$, und die Anzahl der erwarteten Groveriterationen ist kleiner als $m/2$, da $j$ zuf"allig aus $0\le j<m$ ausgew"ahlt wird. 
              Wir sagen, da\3 der Algorithmus die kritische Phase erreicht, wenn er mehr als $\oben{\log_{\lambda} m_0}$ mal die Schleife ausf"uhrt. Die 
              erwartete Anzahl der insgesamt ausgef"uhrten Groveriterationen ist, wenn die kritische Phase erreicht wurde 
              %%\newpage 
              \begin{eqnarray*} 
                \frac{1}{2}\sum_{s=1}^{\ok} \lambda^{s-1} & =   & \frac{1}{2} \sum_{s=0}^{\ok - 1} \lambda^s \\ 
                                                          & =   & \frac{1}{2} \frac{\lambda^{\ok} - 1}{\lambda - 1} \\ 
                                                          & \le & \frac{1}{2} \frac{\lambda^{\ok}}{\lambda - 1} \\ 
                                                          & \le & \frac{1}{2} \frac{\lambda^{\k + 1}}{\lambda - 1} \\ 
                                                          & =   & \frac{1}{2} \frac{\lambda}{\lambda - 1}m_0 \\ 
                                                          & =   & 3 m_0 
              \end{eqnarray*}  
              Falls also der Algorithmus Erfolg hat, bevor er die kritische Phase erreicht, schafft er das in ${\cal O}(m_0) \subset {\cal O}(\sqrt{N/t})$ 
              Schritten. 
              Wenn die kritische Phase erreicht wurde, dann findet der Algorithmus wegen Lemma $1.2.4$ eine L"osung mit einer Wahrscheinlichkeit von mindestens $1/4$, da 
              $m \ge \frac{1}{\sin(2\Theta)}$, bei jedem Durchlauf der Schleife. Die Anzahl der insgesamt ausgef"uhrten Groveriterationen, nachdem die 
              kritische Phase erreicht wurde, wird in der Tabelle angegeben. 
              $$
              \begin{array}{ccc} 
                \mbox{Durchlauf}    & \mbox{Wahrscheinlichkeit}                 & \mbox{ausgef"uhrte Iterationen} \\ 
                                    &                                           &                                 \\
                1                   & \frac{1}{4}                               & \frac{1}{2}\sum_{s=0}^{\ok - 1} \lambda^s + \frac{1}{2}\lambda^{\ok} \\ 
                2                   & (\frac{3}{4})(\frac{1}{4})                & \frac{1}{2}\sum_{s=0}^{\ok + 1} \lambda^s \\ 
                3                   & (\frac{3}{4})^2(\frac{1}{4})              & \frac{1}{2}\sum_{s=0}^{\ok + 2} \lambda^s \\ 
                \vdots              & \vdots                                    & \vdots                                    \\
                j                   & (\frac{3}{4})^{(j-1)}(\frac{1}{4})        & \frac{1}{2}\sum_{s=0}^{\ok + j - 1} \lambda^s
              \end{array}
              $$ 
              Dabei wird durch $1/4$ die Wahrscheinlichkeit, eine L"osung zu finden, nach unten abgesch"atzt und durch $3/4$ die Wahrscheinlichkeit des
              komplment"aren Ereignisses nach oben abgesch"atzt. Der Exponent $j$ bei $(3/4)^j$ gibt also an, wie oft keine L"osung gefunden wurde. Der Summand
              $\frac{1}{2} \sum_{s=0}^{\ok - 1}$ ist die erwartete Anzahl der total ausgef"uhrten Groveriterationen bis die kritische Phase erreicht wurde.
              Damit ergibt sich f"ur die erwartete Anzahl der insgesamt ausgef"uhrten Iterationen nach der kritschen Phase folgende Formel.
              \begin{eqnarray} 
                &     & \sum_{j=1}^{\infty}\frac{1}{4}\left(\frac{3}{4}\right)^{j-1} \left(\frac{1}{2}\sum_{s=0}^{\ok - 1 + j}\lambda^s \right) \nonumber \\ 
                & =   & \frac{1}{8}\sum_{j=0}^{\infty}\left(\frac{3}{4}\right)^j \sum_{s=0}^{\ok + j} \lambda^s \nonumber\\ 
                & =   & \frac{1}{8}\sum_{j=0}^{\infty}\left(\frac{3}{4}\right)^j \frac{\lambda^{\ok + j + 1} - 1}{\lambda - 1} \nonumber\\ 
                & \le & \frac{1}{8}\sum_{j=0}^{\infty}\left(\frac{3}{4}\right)^j \frac{\lambda^{\ok + j + 1}}{\lambda - 1} \nonumber\\ 
                & =   & \frac{1}{8} \frac{\lambda}{\lambda - 1} \lambda^{\ok} \sum_{j=0}^{\infty}\left(\frac{3\lambda}{4}\right)^j\nonumber \\
                & =   & \left(\frac{1}{8}\frac{\lambda}{\lambda - 1} m_0\right) \frac{1}{1-\frac{3\lambda}{4}} \nonumber\\ 
                & =   & \frac{15}{2} m_0 \label{schritte}
              \end{eqnarray} 
              \hfill $\Box$

              Man k"onnte noch besser absch"atzen, da aufgrund der Zuweisung $m=\min(\lambda m, \sqrt{N})$ nie mehr als $\unten{\sqrt{N}}$ Iterationen 
              ausgef"uhrt werden. Diese obere Grenze ist sinnvoll, da bei $t$ L"osungen die Anzahl der notwendigen Iterationen $m\le\frac{\pi}{4}\sqrt{N/t}$
              gilt. $m$ wird bei $t=1$ maximal und es gilt in diesem Fall $m\le\frac{\pi}{4}\sqrt{N} < \sqrt{N}$. Mehr als $\unten{\sqrt{N}}$ Iterationen bringen also keinen 
              Vorteil.
              Dies wird in der obigen Absch"atzung nicht ber"ucksichtigt. Dieser Algorithmus findet also eine L"osung in ${\cal O}( \sqrt{N/t})$ Schritten, 
              wenn $1\le t \le \frac{3}{4}N$ gilt. 
              Der Fall $t=0$ wurde noch nicht zufriedenstellend gel"ost.
              F"ur den Fall $t>\frac{3}{4}N$ erh"alt man klassisch mit einer Wahrscheinlichkeit von mindestens $\frac{3}{4}$ eine L"osung, wenn man ein zuf"alliges $i$ nimmt.

%% ************************** M I N I M U M *********************** 
     \section{Algorithmus zur Bestimmung des Minimums} 
       \subsection{Problemstellung}
         Das Auswahlproblem ist wie folgt definiert: es seien ein unsortiertes Array $V[i]$ mit $i\in X_N=\{0,\ldots,N-1\}$, das $N$ verschiedene Elemente aus einer
         geordneten Menge ${\cal V}$ enth"alt, und eine Zahl $1\le t\le N$ gegeben. Es soll das Element vom Rang t, das hei\3t, dasjenige Element, das gr"o\3er oder gleich als
         $t-1$ Elemente und kleiner als $N-t$ Elemente ist, bestimmt werden. Das Minimum von  $V$ ist das Element vom Rang $1$. Das Minimumsuchproblem besteht darin, den Index $y$ 
         eines Elements zu finden, so da\3 $V[y]$ das Minimum ist. Auf einer klassischen deterministischen oder probabilistischen Turingmaschine ben"otigt dieses Problem 
         ${\cal O}(N)$ Proben. Dieser Algorithmus wurde in \cite{hoyermin} vorgestellt.

         Es wird hier ein Quantenalgorithmus vorgestellt, der das Minimum findet, wobei er lediglich ${\cal O}(\sqrt{N})$ Proben dazu ben"otigt. Die wichtigste
         Subroutine ist der Quantensuchalgorithmus bei unbekannter Anzahl der L"osungen.
       \subsection{Beschreibung des Algorithmus}
         Der Minimumsuchalgorithmus benutzt den Quantensuchalgorithmus, um kleinere Elemente als ein bestimmter Schwellwert $V[y]$, der durch einen Index $y$ bestimmt 
         wird, zu finden 
         Das Ergebnis des Suchalgorithmus wird als neuer Schwellwert gesetzt. Dieser Proze\3 wird so oft ausgef"uhrt, bis die Wahrscheinlichkeit, da\3 der 
         Schwellwert das Minimum ist, ausreichend gro\3 ist.
         \begin{alg} Minimumsuchalgorithmus \end{alg} 
         Die Markierungsfunktion sei
         $$ C(i) = \cases{1 & falls $V[i]<V[y]$ \cr
                          0 & sonst                 }$$
         wobei $y$ der aktuelle Schwellwertindex ist.
         \begin{enumerate} 
           \item W"ahle einen Schwellwertindex $0\le y \le N - 1$ zuf"allig gleichverteilt.
           \item Wiederhole die folgenden Schritte solange die Anzahl der ausgef"uhrten Groveriterationen kleiner als $m_0$ ist.
             \begin{enumerate} 
                \item Setze $m=1$ und $\lambda=\frac{6}{5}$.
                \item W"ahle $j$ zuf"allig gleichverteilt aus den nichtnegativen ganzen Zahlen kleiner als $m$.
                \item F"uhre $j$ Iterationen des Groveralgorithmus beginnend im Zustand $\ket{\Psi_0}=\frac{1}{\sqrt{N}}\sum_{i=0}^{N-1}\ket{i}$ aus.
                \item Messe das Quantenregister. Es sei $i$ das Ergebnis der Messung.
                \item Falls $V[i]<V[y]$, dann setze $y=i$ als neuen Schwellenwertindex und gehe zu (a).
                \item Sonst setze $m=\min(\lambda m,\sqrt{N})$ und gehe zu (b)
             \end{enumerate} 
           \item Messe das Quantenregister. Sei $i$ das Ergebnis. Falls $V[i]<V[y]$ dann setze $y=i$. 
         \end{enumerate} 
         Um die Erfolgswahrscheinlichkeit des Minimumsuchalgorithmus zu bestimmen, soll angenommen werden, da\3 es keine Abbruchbedingung gibt und da\3 der Algorithmus lang
         genug l"auft, um das Minimum zu finden. Dieser Algorithmus wird als Endlos-Algorithmus bezeichnet. Es wird jetzt die mittlere Anzahl der insgesamt ausgef"uhrten
         Groveriterationen, die n"otig sind, um das Minimum zu finden, bestimmt.

         Zu jedem Zeitpunkt sucht der Endlos-Algorithmus das Minimum unter $t$ Elementen, die kleiner sind als $V[y]$. W"ahrend der Ausf"uhrung wird jedes dieser Elemente mit
         gleicher Wahrscheinlichkeit als neuer Schwellwertindex ausgew"ahlt. Das folgende Lemma besagt, da\3 diese Wahrscheinlichkeit der Kehrwert des Ranges des
         Elements ist und nicht von der L"ange des Arrays abh"angt.
         \begin{lemma} Sei $p(t,r)$ die Wahrscheinlichkeit, da\3 der Index eines Elements vom Rang $r$ irgendwann als der Schwellwertindex ausgew"ahlt wird, wenn der
            Endlos-Algorithmus unter $t$ Elementen sucht. Es gilt $p(t,r)=1/r$, falls $r\le t$ und $p(t,r)=0$ sonst.
         \end{lemma}
         \beweis Der Fall $r>t$ ist klar. Der Fall $r\le t$ wird durch vollst"andige Induktion "uber $t$ f"ur ein festes $r$ bewiesen. Der Induktionsanfang $p(r,r)=1/r$ ist
           klar, da der Quantensuchalgorithmus eine beliebige L"osung mit gleicher Wahrscheinlichkeit liefert. Die Induktionsannahme ist, da\3 $p(k,r)=1/r$ f"ur alle $k\in [r,t]$
           gilt. Im Induktionsschritt sind $t+1$ Elemente markiert, und $y$ wird zuf"allig gleichverteilt unter den $t+1$ M"oglichkeiten ausgew"ahlt. Der Algorithmus kann
           ein Element $V[y]$ mit Rang $r$, gr"o\3er als $r$ oder kleiner als $r$ selektieren. Nur die ersten beiden F"alle tragen zu der Summation bei
           \begin{eqnarray*}
              p(t+1,r) & = & \frac{1}{t+1} + \sum_{k=r+1}^{t+1}\frac{1}{t+1} p(k-1,r) \\
                       & = & \frac{1}{t+1} + \frac{t+1-r}{t+1} \cdot \frac{1}{r} \\
                       & = & \frac{1}{r}\left( \frac{r}{t+1} + \frac{t+1-r}{t+1}\right) \\
                       & = & \frac{1}{r}
           \end{eqnarray*}
           \hfill $\Box$

           Die mittlere Anzahl der insgesamt ausgef"uhrten Groveriterationen im Quantensuchalgorithmus betr"agt nach Gleichung \ref{schritte} h"ochstens $\frac{15}{2}\sqrt{N/t}$, 
           falls $t$ Elemente unter
           $N$ markiert sind. Damit kann man die obere Schranke f"ur die mittlere Anzahl der total ausgef"uhrten Groveriterationen $m_0$, die notwendig sind, um das 
           Minimum zu finden, ausrechnen.
           \begin{eqnarray*}
              \sum_{r=2}^{N}p(N,r)\frac{15}{2} \sqrt{\frac{N}{r-1}} & =   & \frac{15}{2}\sqrt{N}\sum_{r=1}^{N-1}\frac{1}{r+1} \frac{1}{\sqrt{N}} \\
                                                                    & \le & \frac{15}{2}\sqrt{N}\left(\frac{1}{2} + \sum_{r=2}^{N-1}r^{-3/2}\right) \\
                                                                    & \le & \frac{15}{2}\sqrt{N}\left(\frac{1}{2} + \int_{r=1}^{N-1}r^{-3/2} dr\right) \\
                                                                    & =   & \frac{15}{2}\sqrt{N}\left(\frac{1}{2} + \left[-2r^{-1/2}\right]_{r=1}^{N-1}\right) \\
                                                                    & =   & \frac{15}{2}\sqrt{N}\left(\frac{1}{2} - \frac{2}{\sqrt{N-1}} + 2\right) \\
                                                                    & \le & \frac{95}{4}\sqrt{N}
           \end{eqnarray*}
           \hfill $\Box$
           \begin{satz}
             Der Minimumsuchalgorithmus findet den Index des Minimums mit einer Wahrscheinlichkeit von mindestens $1/2$.
           \end{satz}
           \beweis 
             Sei $p(t)$ die Wahrscheinlichkeit, da\3 das Minimum nach genau $t$ Groveriterationen gefunden wird. Annahme $p_0 = p(2m_0) < 1/2$. Damit ergibt sich f"ur die
             mittlere Anzahl der total ausgef"uhrten Groveriterationen $> p_0 \cdot 0 + (1-p_0)2m_0 > m_0$. Das ist ein Widerspruch. Nach $2m_0$ Groveriterationen ist die
             Erfolgswahrscheinlichkeit also gr"o\3er als $1/2$. \hfill $\Box$

           Die Erfolgswahrscheinlichkeit kann vergr"o\3ert werden, indem man den Algorithmus $c$-mal ausf"uhrt. Sei $y$ so, da\3 $V[y]=\min\{V[y_1],\ldots,V[y_c]\}$ gilt.
           Dann ist $V[y]$ das Minimum des Arrays $V$ mit einer Wahrscheinlichkeit von mindestens $1-1/2^c$. Die Erfolgswahrscheinlichkeit ist sogar besser, wenn man den 
           Algorithmus mit der Abbruchbedingung $c2m_0$ ausf"uhrt, da jedesmal die Information aus den vorhergehenden Suchergebnissen ausgenutzt wird.
              
           


%% ************************** M E D I A N ************************* 
     \section{Algorithmus zur Sch"atzung des Medians}
       \subsection{Problemstellung}
       Das Auswahlproblem ist wie folgt definiert: es seien ein unsortiertes Array $V[i]$ mit $i\in X_N=\{0,\ldots, N-1\}$, das $N$ verschiedene Elemente 
       aus einer geordneten Menge ${\cal V}$ enth"alt, und eine Zahl $1\le t\le N$ gegeben. Es soll
       das $t$-gr"o\3te Element in $V$, d.h. dasjenige Element, das gr"o\3er als $t-1$ Elemente und kleiner als $N-t$ Elemente ist, bestimmt werden. Der Median
       von $V$ ist das $\oben{N/2}$-gr"o\3te Element. 
              
       Betrachten wir zun"achst ein einfacheres Problem, auf das sich das Problem, den Median zu bestimmen, zur"uck\-f"uhren l"a\3t. Das Array $V$ werde durch einen 
       Schwellenwertindex $\mu\in{\cal V}$ in zwei Teilmengen $A=\{V[i] | V[i]<\mu\}$ und in
       $B=\{V[i] | V[i]\ge\mu\}$ partitioniert, wobei die M"achtigkeiten der beiden Mengen $\betrag{A}=(N/2)(1+\epsilon)$ bzw.
       $\betrag{B}=(N/2)(1-\epsilon)$ seien. Dabei gelte $0.1\epsilon_0 < \betrag{\epsilon} < \epsilon_0$ f"ur ein festes $\epsilon_0$, das kleiner als $0.1$ ist.
       Die Aufgabe ist nun, $\betrag{\epsilon}$ mit einer Genauigkeit $\Delta$ zu bestimmen. Wenn $\widetilde\epsilon$ ein Sch"atzwert f"ur $\betrag{\epsilon}$ ist, dann soll 
       f"ur $\widetilde\epsilon$ mit einer hohen Wahrscheinlichkeit $\widetilde{\epsilon}(1-\Delta) < \betrag{\epsilon} < \widetilde{\epsilon}(1+\Delta)$ gelten.
       Dieses Problem l"a\3t sich auf einem Quantencomputer l"osen, indem man die Anzahl der Zust"ande $\ket{i}$ mit $C(i)=0$ sch"atzt, wobei
       $$ C(i) = \cases{ 1 & falls $V[i] \ge\mu$ \cr
                         0 & falls $V[i] < \mu$ }. $$
       Grovers Algorithmus ben"otigt dazu ${\cal O}(\frac{1}{\epsilon_0 \Delta^2})$ Proben. Klassische Algorithmen ben"otigen 
       $\Omega(\frac{1}{\epsilon_0^2\Delta^2})$.

       \subsection{Beschreibung des Algorithmus}
         Die Amplituden der Zust"ande $\ket{i}$ mit $C(i)=0$ werden als $k$-Amplituden und die Amplituden der Zust"ande $\ket{j}$ mit $C(j)=1$ werden als
         $l$-Amplituden bezeichnet. \\
         \noindent
         Eine bedingte Phasendrehung $S_{F,\delta}$ operiert wie folgt auf den Zust"anden $\ket{i}$ mit $i\in X_N$
         $$ S_{F,\delta} \ket{i} = \left\{
                                     \begin{array}{rc}
                                       e^{i \delta}\ket{i} & \mbox{ falls } F\left( i\right) = 1 \\
                                                   \ket{i} & \mbox{ sonst, }
                                     \end{array}
                                   \right.
         $$
         wobei die Funktion $F:X_N \to \{0,1\}$ eine Bedingung beschreibt und der Winkel $\delta\in (0,2\pi )$ angibt, um wieviel die Phase gedreht werden soll. 
         Es werden noch folgende Transformationen ben"otigt
         \begin{enumerate}
           \item Diffusionstransformation $D=-I+2P$,
           \item Schiebetransformationen $S=-iI + (1+i)P$.
         \end{enumerate}
         \begin{lemma}[Schiebetransformation]
           Die Schiebetransformation l"a\3t als $S = W\diag{1,-i,\ldots,-i} W^{\dagger}$ schreiben.
         \end{lemma}
         \beweis Analog zu Satz~\ref{diffusion} \hfill $\Box$
         \begin{alg}[Bestimmung von $\epsilon$] Es soll die Abweichung $\epsilon$ zwischen den Zust"anden $\ket{i}$ mit $C(i)=0$ und $\ket{j}$ mit $C(j)=1$ 
           gesch"atzt werden.
         \end{alg}
         \begin{enumerate}
         \item $\nu={\cal O}(\frac{1}{\Delta^2})$ Messwerte ermitteln
         \begin{enumerate}
           \item Gleichverteilung $ \ket{\Psi_0} = \sum_{i=0}^{N-1} \ket{i} $
           \item bedingte Phasendrehung $S_{C,\frac{\pi}{2}}$ 
           \item Schiebetransformation $S$ 
           \item wiederhole die folgenden Schritte $m={\cal O}(\frac{1}{\epsilon_0})$ mal
             \begin{enumerate}
               \item bedingte Phasendrehung $S_{\neg C,\pi}$ 
               \item Diffusionstransformation $D$ 
               \item bedingte Phasendrehung $S_{C,\pi}$
               \item Diffusionstransformation $D$ 
             \end{enumerate}
           \item Messung des Registers und Speicherung des Messwerts
         \end{enumerate}
         \item Sch"atzung von $\betrag{\epsilon}$ durch statistische Auswertung der Messwerte
         \end{enumerate}
         %% **************** Funktionsweise des Algorithmus ****************
         \subsection{Funktionsweise des Algorithmus}
         \subsubsection*{1. Teil : Messwerte ermitteln}
           $$ 
           \begin{array}{rlcc}
             \mbox{1.(a)} & \mbox{Herstellung der Gleichverteilung }            & k=1         & l=1 \cr
             \mbox{1.(b)} & \mbox{bedingte Phasendrehung } S_{C,\frac{\pi}{2}}  & k=1         & l=i \cr
             \mbox{1.(c)} & \mbox{Schiebetransformation } S=-iI+(1+i)P          & k=\epsilon  & l=(1+\epsilon)+i
           \end{array}
           $$
           \begin{lemma}
             Anwendung der Schiebetransformation auf $k=1$ und $l=i$ bewirkt $k=\epsilon$ und $l=(1+\epsilon)+i$.
           \end{lemma}
           \beweis F"ur den Durchschnittswert $\alpha$ gilt
             \begin{eqnarray*}
               \alpha & = & \frac{1}{N}\left( \frac{N}{2}\left(1+\epsilon\right)+\frac{N}{2}\left(1-\epsilon\right)i\right) \\
                      & = & \frac{1}{2}\left(1 + \epsilon + i - i\epsilon\right)
             \end{eqnarray*}
             Aus der Zerlegung $S=-iI+(1+i)P$ ergibt sich nun f"ur die $k$-Amplituden
             \begin{eqnarray*}
               k & = & -i + \frac{1}{2}\left(1+i\right)\left(1+\epsilon + i - i\epsilon \right) \\
                 & = & -i + \frac{1}{2}\left(1+ \epsilon + i - i\epsilon + i + i\epsilon - 1 + \epsilon\right) \\
                 & = & -i + \frac{1}{2}\left(2\epsilon + 2i\right) \\
                 & = & \epsilon .
             \end{eqnarray*}
             Die $l$-Amplitude $l = \left(1+\epsilon\right) + i$ kann analog berechnet werden. \hfill $\Box$
           
             Die Schiebetransformation $S$ verlagert die Amplituden zu den Zust"anden $\ket{j}$ mit $C(j)=1$, das hei\3t, die $k$-Amplituden werden ungef"ahr gleich Null und
             der Betrag der $l$-Amplituden wird ungef"ahr $\sqrt{2}$ gro\3 sein. Falls die Anzahl der Zust"ande $\ket{i}$ mit $C(i)=0$ gleich der Anzahl der 
             Zust"ande $\ket{j}$ mit $C(j)=1$ ist, dann sind nach der Schiebetransformation die $k$-Amplituden genau gleich Null und die $l$-Amplituden gleich $1+i$.  \\
             \noindent
             Das Herz des Algorithmus ist die Iteration im Schritt 1.(d), die die Betr"age der $k$-Amplituden linear verst"arkt. Um diesen Verst"arkungsproze\3 zu verstehen, 
             mu\3 man wissen, wie die Diffusionstransformation $D$ die $k$- und $l$-Amplituden
             ver"andert. Zuerst wird die Wirkung der Diffusionstransformation bei zwei einfachen Konfigurationen der $k$- und $l$-Amplituden untersucht. Der 
             allgemeine Fall ergibt sich aus der linearen Superposition der beiden Konfigurationen.
             \begin{lemma}
               Anwendung der Diffusionstransformation $D$ auf $k=1$ und $l=0$ bewirkt $k=\epsilon$ und $l=1+\epsilon$.
             \end{lemma}
             \beweis
               Der Durchschnittswert ist $\alpha = \frac{1}{2}(1+\epsilon)$.
               $$
                 \begin{array}{ccccc}
                   k & = & -1 + 2 \alpha  & = & \epsilon \\ 
                   l & = & 0 + 2 \alpha   & = & 1 + \epsilon 
                 \end{array}
               $$  \hfill $\Box$
             \begin{lemma}
               Anwendung der Diffusionstransformation $D$ auf $k=0$ und $l=1$ bewirkt $k=1-\epsilon$ und $l=-\epsilon$.
             \end{lemma}
             \beweis 
               Der Durchschnittswert ist $\alpha = \frac{1}{2}(1-\epsilon)$.
               $$
                 \begin{array}{ccccc}
                    k & = & 0 + (1-\epsilon)  & = & 1 - \epsilon \\ 
                    l & = & -1 + (1-\epsilon) & = & -\epsilon
                 \end{array}
                $$ \hfill $\Box$

             \begin{lemma} Schritt 1.(d) bei Anfangskonfiguration $k=1$ und $l=0$
             \end{lemma}
             \begin{tabular}{rlll}
               i.    & bedingte Phasendrehung $S_{\neg C,\pi}$ & $k = -1$        & $l = 0$                 \\
               ii.   & Diffusionstransformation $D$            & $k = -\epsilon$ & $l = -(1+\epsilon)$     \\
               iii.  & bedingte Phasendrehung $S_{C,\pi}$      & $k = -\epsilon$ & $l = 1+\epsilon$        \\
               iv.   & Diffusionstransformation                &                 &
             \end{tabular}
             $$ 
             \begin{array}{rclrcl}
               k & = & -\epsilon \cdot \epsilon + (1 + \epsilon)(1 - \epsilon)  & l & = & -\epsilon (1 + \epsilon) + (1 + \epsilon)(-\epsilon) \\
                 & = & -\epsilon^2 + 1 -\epsilon^2                              &   & = & -\epsilon - \epsilon^2 - \epsilon - \epsilon^2 \\
                 & = & 1 - 2\epsilon^2                                          &   & = & -2\epsilon - 2\epsilon^2 
             \end{array}
             $$ \hfill $\Box$
             \begin{lemma} Schritt 1.(d) bei Anfangskonfiguration $k=0$ und $l=1$
             \end{lemma}
             \begin{tabular}{rlll}
               i.   & bedingte Phasendrehung $S_{\neg C,\pi}$ & $k = 0 $         & $l = 1$ \\
               ii.  & Diffusionstransformation $D$            & $k = 1-\epsilon$ & $l = -\epsilon$ \\
               iii. & bedingte Phasendrehung $S_{C,\pi}$      & $k = 1-\epsilon$ & $l = \epsilon$ \\
               iv.  & Diffusionstransformation                &                  &
             \end{tabular}
             $$
             \begin{array}{rclrcl}
               k & = & (1-\epsilon)\epsilon + \epsilon(1-\epsilon)       & l & = & (1 -\epsilon )(1 + \epsilon) + \epsilon(-\epsilon) \\
                 & = & \epsilon - \epsilon^2 + \epsilon - \epsilon^2     &   & = & 1 - \epsilon^2 - \epsilon^2 \\
                 & = & 2\epsilon - 2 \epsilon^2                          &   & = & 1 -2\epsilon^2
             \end{array}
             $$ \hfill $\Box$
             \begin{satz}[Geschlossene Formeln f"ur die $k$- und $l$-Amplituden]
             Mit Hilfe der beiden vorhergehenden Lemmata l"a\3t sich die Wirkung eines Iterationsschrittes im Schritt  1.(d) angeben:
             \begin{eqnarray*}
               k_j & = & \epsilon\cos(j\Theta) + \gamma(1+\epsilon+i)\sin(j\Theta) \\
               l_j & = & -\frac{\epsilon}{\gamma}\sin{j\Theta}+(1+\epsilon+i)\cos(j\Theta),
             \end{eqnarray*}
             wobei $\cos(\Theta) = 1 - 2\epsilon^2$ und $\gamma^2 = \frac{2\epsilon-2\epsilon^2}{2\epsilon+2\epsilon^2}$. 
             \end{satz}
             \beweis Es gilt 
             $$ \pmatrix{k_{j+1} \cr l_{j+1}} = \pmatrix{ 1 - 2\epsilon^2           &  2\epsilon - 2\epsilon^2 \cr
                                                          -2\epsilon - 2\epsilon^2  &  1 - 2\epsilon^2} 
                                                \pmatrix{k_j \cr \l_j} $$
             mit den Anfangsbedingungen $k_0 = \epsilon$ und $l_0 = (1+\epsilon)+i$. \\
             \noindent
             Mit den Substitutionen $\cos(\Theta) = 1 - 2\epsilon^2$ und $\gamma^2 = \frac{2\epsilon-2\epsilon^2}{2\epsilon+2\epsilon^2}$ erh"alt man
             $$ \pmatrix{k_{j+1} \cr l_{j+1}} = \pmatrix{ \cos\Theta                    &  \gamma\sin\Theta \cr
                                                          -\frac{1}{\gamma}\sin\Theta   &  \cos\Theta }
                                                \pmatrix{k_j \cr \l_j} $$
             Die Potenzen dieser Matrix lassen sich sehr einfach berechnen. Mit vollst"andiger Induktion beweist man leicht, da\3
             $$ \pmatrix{k_j \cr l_j} = \pmatrix{ \cos(j\Theta)                    &  \gamma\sin(j\Theta) \cr
                                                  -\frac{1}{\gamma}\sin(j\Theta)   &  \cos{j\Theta} }
                                                \pmatrix{\epsilon \cr (1+\epsilon)+i} $$
             gilt.
             Damit lassen sich folgende geschlossene Formeln f"ur $k_j$ und $l_j$
             \begin{eqnarray*}
               k_j & = & \epsilon\cos(j\Theta) + \gamma(1+\epsilon+i)\sin(j\Theta) \\
               l_j & = & -\frac{\epsilon}{\gamma}\sin(j\Theta)+(1+\epsilon+i)\cos(j\Theta)
             \end{eqnarray*}
             ablesen. \hfill $\Box$
             \begin{lemma}
               Es existiert eine Zahl $m < \frac{1}{40\epsilon_0}$, so da\3 nach $m$ Iterationen f"ur den Betrag der $k$-Amplituden 
               $\betrag{k_m} > \frac{\betrag{\epsilon}}{20\epsilon_0}$ gilt.
             \end{lemma}    
             \beweis Es werden Ausdr"ucke mit trigonometrischen Funktionen abgesch"atzt, wobei die Tatsache ausgenutzt wird, da\3 $\betrag{\epsilon}$ klein ist. 
             Aufgrund der Definition gilt $\cos\Theta=1 - 2\epsilon^2$. Mit den folgenden Identit"aten 
             $$
               \cos^2\Theta = \left( 1 - 2\epsilon \right)^2 = 1-4\epsilon^2+4\epsilon^4 \\
             $$
             $$
               \sin^2\Theta = 1 - \cos^2\Theta = 1 - 1 + 4\epsilon^2 - 4\epsilon^4 = 4\epsilon^2 - 4\epsilon^4 
             $$
             $$
                 \sin\Theta = \sqrt{4\epsilon^2 - 4\epsilon^4} = 2\epsilon\sqrt{1-\epsilon^2}
             $$ 
             kann man folgende Absch"atzungen $\sin\Theta \approx 2\epsilon$ und $\Theta \approx \sin\Theta$ machen, da die Sinus-Funktion f"ur kleine Argumente linear ist.
             Es gilt $\gamma \approx 1$. F"ur nicht zu grosse $j$ gilt $\sin(j\Theta) \approx j\Theta$ und $\cos(j\Theta) \approx j\Theta$. 
             Mit diesen Absch"atzungen erh"alt man f"ur die $k$-Amplitude
             \begin{eqnarray*}
                      k_j   & =       & \epsilon\cos(j\Theta)+\gamma\left(1+\epsilon+i\right)\sin(j\Theta) \\
                            & \approx & \epsilon + (1 + \epsilon + i)2j\epsilon
             \end{eqnarray*}
             und den Betrag der $k$-Amplitude
             $$ \betrag{k_j} \approx 2\sqrt{2}\betrag{\epsilon}j, $$
             indem man alle Summanden mit $\epsilon^3, \epsilon^4$ und $j\epsilon^2$ wegl"a\3t.
             F"ur kleine $\epsilon$ h"angt also der Betrag der $k$-Amplitude linear von $\betrag{\epsilon}$ ab. 
             Mit dieser linearen Approximation sieht man, da\3 eine Zahl $m$ existiert, so da\3 $m$ Iterationen f"ur den Betrag der $k$-Amplitude $\betrag{k_m}$
             $$ \betrag{k_m} > \frac{\betrag{\epsilon}}{20\epsilon_0} $$ gilt.
             Es mu\3 f"ur $m$ 
             \begin{eqnarray*}
               2\sqrt{2}m\betrag{\epsilon} & > & \frac{\betrag{\epsilon}}{20\epsilon_0} \\
                                        m  & > & \frac{1}{40\sqrt{2} \epsilon_0}
             \end{eqnarray*} 
             gelten. Man sieht insbesondere, da\3 die Anzahl der Iterationen kleiner als $m < 1/40\epsilon_0$ gew"ahlt werden kann.
             Dies ist auch der Grund, warum Grovers Algorithmus nur ${\cal O}(1/\epsilon_0 \Delta^2)$ 
             Proben braucht, um $\betrag{\epsilon}$ zu bestimmen, wohingegen ein klassischer Algorithmus $\Omega(1/\epsilon_0^2\Delta^2)$
             ben"otigt. \hfill $\Box$
             
             Das folgende Mathematica-Programm zeichet den Betrag der $k$-Amplitude und ihre lineare Approximation.
             \begin{center}
             {\footnotesize
             \begin{verbatim}
(* Verst"arkungsprozess in 1.(d), PW                          *)
(* zeichnet den Betrag der k-Amplitude und ihre Approximation *)
(* in Abh"angigkeit der Anzahl der Iterationen                *)

epsilon = 0.1;

Theta = ArcCos[1 - 2 epsilon^2];
gamma = Sqrt[(2 epsilon - 2 epsilon^2) / (2 epsilon + 2 epsilon^2)]; 

Plot[ {2 Sqrt[2] Abs[epsilon] j, Abs[ epsilon Cos[j Theta] + gamma ( 1 + epsilon + I) Sin[j Theta]]}, {j,0,50}] 
             \end{verbatim}
             }
             \end{center}           
             Man sieht mit diesem Programm, da\3 die Approximation sehr gut ist.
             %%
             %% 
             \subsection*{2. Teil: Sch"atzung von $\betrag{\epsilon}$}
             %%
             %%
             Im 2.Teil werden die Messwerte statistisch ausgewertet, um $\betrag{\epsilon}$ zu sch"atzen. Die Grundlage dieser Auswertung ist das 
             Gesetz der gro\3en Zahlen ben"otigt.
             \begin{satz}[Gesetz der gro\3en Zahlen]
               Es sei in einem Bernoulli-Experiment $p$ die Wahrscheinlichkeit eines Ereignisses $A$ und $(1-p)$ die Wahrscheinlichkeit des komplement"aren 
               Ereignisses $\overline{A}$. Bei $\nu$ Wiederholungen liegt $f$, die relative H"aufigkeit des Ereignisses $A$, zwischen $p-\kappa\sqrt{1/\nu}$ und
               $p+\kappa\sqrt{1/\nu}$ mit einer Wahrscheinlichkeit von mindestens $(1-2\exp(-2\kappa^2))$. \hfill $\Box$
             \end{satz}
             %
             \begin{lemma}[Genauigkeit von $\epsilon$] Das Experiment im ersten Teil werde $\nu = {\cal O}(1/\Delta^2)$ mal ausgef"uhrt. $f$ sei der Bruchteil der Zust"ande 
               $\ket{i}$ mit
               $C(i)=0$, die im Schritt 1.(e) gemessen wurden. Mit einer Wahrscheinlichkeit von mindestens $(1-2\exp(-2\kappa^2))$ liegt die Gr"o\3e 
               $ \frac{1}{2}(\betrag{k_m }^2(1+\epsilon))$ in dem Intervall von $(f-\kappa\Delta)$ bis $(f+\kappa\Delta)$.
             \end{lemma}
             \beweis Die Wahrscheinlichkeit, irgend einen Zustand $\ket{i}$ mit $C(i)=0$ 
             zu messen, betr"agt $\frac{1}{2}(\betrag{k_m }^2(1+\epsilon))$, da es $(N/2)(1+\epsilon)$ solcher Zust"ande gibt. 
          
             Das Gesetz der gro\3en Zahlen besagt, da\3 nach $\nu$ Wiederholungen $f$ in dem Intervall von $\frac{1}{2}(\betrag{k_m }^2(1+\epsilon))-\kappa\Delta$ bis
             $\frac{1}{2}(\betrag{k_m }^2(1+\epsilon))+\kappa\Delta$ mit einer Wahrscheinlichkeit von mindestens $(1-2\exp(-2\kappa^2))$ liegt. Durch Umformen 
             sieht man, da\3 $\frac{1}{2}(\betrag{k_m }^2(1+\epsilon))$ zwischen $f-\kappa\Delta$ und $f+\kappa\Delta$ mit einer Wahrscheinlichkeit von mindestens
             $(1-2\exp(-2\kappa^2))$ liegt. \hfill $\Box$
             %
             \begin{lemma} Seien $a,b\in\R$ mit $a<b$ und $\Delta f = f(b)-f(a)$. Es gelte f"ur die Ableitung der Funktion $f'(x)>d$ f"ur alle $x\in (a,b)$ mit
                $d$ positiv. Dann gilt $b-a<\Delta f / d$.
             \end{lemma}
             \beweis Aufgrund des Zwischenwertsatzes gilt $f(b)-f(a)=f'(\xi)(b-a)$ f"ur ein $\xi\in (a,b)$. Damit folgt $b-a=\Delta f / f'(\xi) < \Delta f / d$.
               \hfill $\Box$
             %
             \begin{satz}
               Der Wert von $\betrag{\epsilon}$ kann mit einem maximal erwarteten Fehler von ${\cal O}(\Delta\epsilon)$ aus dem Sch"atzwert $f$ f"ur
               $f(\epsilon)=\frac{1}{2}(\betrag{k_m }^2(1+\epsilon))$ gewonnen werden.
             \end{satz}
             \beweis Der Sch"atzwert $f$ liegt mit einer Wahrscheinlichkeit von mindestens $(1-\exp(-2\kappa^2))$ in dem Intervall von 
               $f(\epsilon)-\kappa\Delta$ bis $f(\epsilon)+\kappa\Delta$. Sei also $\betrag{f(a)-f(b)}=2\kappa\Delta$.
               \begin{enumerate}
                 \item $\betrag{k_m(\epsilon)}^2 = (\gamma(1 + \epsilon)\sin(m\Theta) + \epsilon\cos(m\Theta))^2 + \gamma^2\sin^2(m\Theta) $
                 \item Um aus der Sch"atzung $f$ f"ur $f(\epsilon)$ den Wert $\betrag{\epsilon}$ sch"atzen zu k"onnen, mu\3 man die partielle Ableitung der Funktion $f(\epsilon)$
                   nach $\epsilon$ kennen. F"ur sie gilt
                   \begin{eqnarray*}
                     {\partial\over\partial\epsilon}f(\epsilon) & =       & {\partial\over\partial\epsilon}\frac{1}{2}(\betrag{k_m}^2(1+\epsilon)) \\
                                                                & \approx & \frac{1}{2}{\partial\over\partial\epsilon}(2\sqrt{2}m\betrag{\epsilon})^2 \\
                                                                & =       & 8m^2\epsilon.
                   \end{eqnarray*}
                 \item Aufgrund des vorhergehenden Lemmas gilt $\betrag{b-a} < {\betrag{f(b)-f(a)}}/{\min\betrag{f'(x)}}$ mit $x\in (a,b)$. Damit und mit
                       $m > {1}/{20\epsilon_0}$ und $\betrag{\epsilon_0 / \epsilon}<10$ erh"alt man 
                       folgende Absch"atzung
                       $$ \betrag{b-a} < \frac{2\kappa\Theta}{8m^2\epsilon} = \frac{\kappa\Theta}{4m^2\epsilon_0^2}\left(\frac{\epsilon_0^2}{\epsilon^2}\right)\betrag{\epsilon} <
                          \frac{\kappa\Theta}{4\left(\frac{1}{20}\right)^2}(10)^2\betrag{\epsilon} = 10000\kappa\Delta\betrag{\epsilon} $$
                       Der Wert von $\betrag{\epsilon}$ kann also aus der relativen H"aufigkeit der Zust"ande $\ket{i}$ mit $C(i)=0$ gesch"atzt werden. Der Sch"atzwert liegt dabei
                       mit einer Wahrscheinlichkeit von mindesten $(1-\exp(-2\kappa^2))$ bei einer Abweichung von h"ochstens $10000\kappa\Delta\epsilon_0$ vom 
                       tats"achlichen Wert von $\betrag{\epsilon}$ vom entfernt. \hfill $\Box$
               \end{enumerate}
               Indem man den Schwellwert $\mu$ leicht varriert, kann man das Vorzeichen von $\epsilon$ feststellen. Man wei\3 jetzt, ob der Schwellwert gr"o\3er oder
               kleiner ist als der Median. Der Schwellwert wird solange ver"andert, bis $\betrag{\epsilon}$ klein genug ist. Dies kann mit dem folgenden
               Programm erreicht werden, das die Quantensubroutine, die $\betrag{\epsilon}$ bestimmt, aufruft, um den Schwellenwert $\mu$ geeignet zu "andern. F"ur den
               Startwert mu\3 $\betrag{\epsilon} < \epsilon_0$ gelten.
\newpage
\begin{verbatim}
median(min, max)
{
  lower = min;
  upper = max;
  while (upper - lower > 1)
  {
    mu = (upper + lower) / 2;
    // liefert epsilon mit Vorzeichen
    // epsilon(schwellwert, epsilon_0, delta)
    est = epsilon_est(mu, 0.1, 0.1); 
    if (est > 0) upper = mu;
    else         lower = mu;
  }
  return mu;
}
\end{verbatim}

               \begin{thebibliography}{10}
                 \bibitem[Ben94]{bernstein} 
                  Bennett, Charles~H. and Bernstein, Ethan and 
                  Brassard, Gilles and Vazirani, Umesh, 
                  {\em Strengths and weaknesses of quantum computing}, LANL preprint quant--ph/9701001
                 \bibitem[Gro96a]{grosuch}
                   Lov K. Grover, {\em A fast quantum mechanical algorithm for database search}, 
                   LANL preprint quant--ph/9605043
                 \bibitem[Gro96b]{grovermedian}
                   Lov K. Grover, {\em A fast quantum mechanical algorithm for estimating the median}, 
                   LANL preprint quant--ph/960702
                 \bibitem[Boy96]{hoyer}
                   Boyer, Michel and Brassard, Gilles and H{\o}yer,
                   Peter and Tapp, Alain,
                   {\em Tight bounds on quantum searching}, LANL preprint quant--ph/9605034
                 \bibitem[Dur96]{hoyermin}
                   D{\"u}rr, Christoph and H{\o}yer, Peter,
                   {\em A quantum algorithm for finding the minimum}, LANL preprint quant--ph/9607014
                 \bibitem[Kit96]{kitaev}
                   A. Yu. Kitaev, {\em Quantum measurements and the Abelian Stabilizer Problem}, LANL preprint quant--ph/9511026
               \end{thebibliography}

 
 
 
 
 
 
 
 
 
 
 
 
 

